\section{Step 1: Build a very simple batched IVP solver}

\begin{frame}{The simplest case}
    
\begin{itemize}
\item Explicit time-stepping (forward Euler), constant time step.
\item Use Python \texttt{CuPy} interface, with \textbf{Raw Kernels} for solvers.
\item To include custom RHS, embed text directly into the CUDA kernel (written in C) before compiling.
\begin{itemize}
    \item Best of both worlds: \texttt{Python} for user interface, low-level \texttt{C} makes code very efficient.
\end{itemize}
\item Compare two different kernels:
\begin{itemize}
    \item \texttt{forward\_euler\_kernel}: on one thread, takes a single time step. Measure time for many steps.
    \item \texttt{forward\_euler\_kernel\_multistep}: on one thread, runs a single IVP to completion. Measure time for a single launch
\end{itemize}
\end{itemize}
\end{frame}

\begin{frame}{Single Time Step vs. All at Once}
% \begin{center}
%     \includegraphics[width=0.4\linewidth]{../figs/forward_euler_singlestep.png}
%     \includegraphics[width=0.4\linewidth]{../figs/forward_euler_multistep.png}
% \end{center}
We can see that it is much faster to do time-stepping \textbf{inside the kernel}.
\end{frame}


\begin{frame}{Comparison to Serial Algorithm}
    Use numpy to solve Lorenz system.
\end{frame}

\begin{frame}{RK23 Implementation}
\begin{itemize}
\item Same thing as before: one ODE per thread.
\item Introduce thread divergence when we decide to accept/reject time step. \textit{Ignore this for now.}
\item Otherwise each thread does the exact same thing. Some threads may finish before others, but this is fine.
\end{itemize}
\end{frame}

\begin{frame}

\end{frame}